\subsection{The specified transformed mechanism}\label{sec:v5-primal}
We apply Proposition~\ref{prop:reserve-preservation} to the complete seed
of Section~\ref{sec:refined-mechanism}.
\begin{lemma}[Zero-value exclusion]\label{lem:zero-item}
At every profile $z$, the selected allocation of $M_*$ satisfies
$z_{ij}=0\Rightarrow x^*_{ij}(z)=0$, including all ties.
\end{lemma}
\begin{proof}
In a base menu each singleton costs at least $1/2$, and each bundle
upgrade costs at least $c=157/500$. In a free $Q$ menu both upgrades
cost at least $b_0-A>0$. In a constrained $Q$ menu, $k\ge c$ and
$C\ge5/6+3k^2/4$, with safe price $B=\min(A,C-k)$.
These imply positive singleton prices and positive bundle upgrades on
the complete certified region. The uniform safe-price and safe-upgrade
lower bounds can be taken as $525947/10^6$ and $721841/3000000$.
In an $E$ lottery, positive safe allocation requires safe value at least
$c$, and positive scarce allocation requires scarce value at least
$A-q-\delta(A)\ge1364159/3000000>0$.
Thus any row allocating a zero-valued item is strictly dominated by an
available row deleting that item. The empty row is selected whenever
maximum utility is zero; positive ties remain between maximizing rows
that have the same exclusion property. The base, $Q$ and $E$ cases cover
the complete mechanism.
\end{proof}


\begin{proposition}[Complete specified mechanism]\label{prop:v5-feasible}
Use $\tau=83/10000$ in \eqref{eq:reserve-transform}, with the seed menus
and joint tie rule of Section~\ref{sec:refined-mechanism}. This defines
a Borel, pointwise DSIC, ex-post IR and jointly feasible mechanism.
\end{proposition}
\begin{proof}
The seed satisfies Proposition~\ref{prop:refined-feasible} and
Lemma~\ref{lem:zero-item}. Its nonnegative payments are at most two by IR.
Proposition~\ref{prop:reserve-preservation} therefore applies at every report.
\end{proof}

This change need not retain a subset of the original allocation at the
same report. Throughout the closed box of halfwidth $10^{-4}$ centered at
\[
 (v_{11},v_{12},v_{21},v_{22})=(1577/2000,123/200,3/5,4/5),
\]
$M_*$ gives both items to bidder 1, whereas $M_\tau$ gives item 1 to
bidder 1 and item 2 to bidder 2. Exact branch inequalities first place
the entire box in the stated base-menu regime. All remaining utility
differences are affine there; their vertex bounds are at least $1/1250$
before the change and $1719/5000000$ after it. The box has volume
$1/625000000000000$. This is a jointly feasible ownership transfer,
not a sampled allocation comparison.

\subsection{Exact revenue through face moments}
For $S\subseteq\{1,2,3,4\}$ let $R_S$ and $N_S$ be the integrals of
the total payment and total expected number of allocated items of $M_*$
on the face where coordinates in $S$ are zero, with all remaining
coordinates independent uniform. Each transformed coordinate has law
$\tau\delta_0+s_\tau\operatorname{Unif}[0,1]$. Independence gives the exact
degree-five polynomial
\begin{equation}\label{eq:v5-revenue}
 R_\tau=R(M_\tau)=
 \sum_{S\subseteq\{1,2,3,4\}}\tau^{|S|}s_\tau^{4-|S|}
                 (s_\tau R_S+\tau N_S),\qquad \tau=83/10000.
\end{equation}
This formula, together with the complete menus, is an exact revenue
characterization. The numerical value is certified by an interval;
neither its midpoint nor an endpoint is asserted to be the exact revenue.
Boundary faces cannot be dropped: they have positive probability after
transformation, even though they had zero four-dimensional volume before it.

For clarity, the sums over faces of each codimension are
\[
\begin{split}
 (\mathcal R_0,\ldots,\mathcal R_4)
   &=(R_*,4R_3,2(R_{\rm SB}+R_{\rm same}+R_{\rm cross}),4R_1,0),\\
 (\mathcal N_0,\ldots,\mathcal N_4)
   &=(N_4,4N_3,2(N_{\rm SB}+N_{\rm same}+N_{\rm cross}),4N_1,0).
\end{split}
\]
Here $R_3,N_3$ refer to one fixed zero coordinate; $R_1,N_1$ to one
remaining nonzero coordinate. Of the two-dimensional faces, one bidder
alone yields $R_{\rm SB}=4/9+2\sqrt2/27$ and
$N_{\rm SB}=(8+2\sqrt2)/9$; the same-item face yields
$R_{\rm same}=31/81$, $N_{\rm same}=5/9$; and the cross-item face
retains the two opposing singleton menus. Also $R_1=2/9$, $N_1=1/3$.
The remaining coefficients reduce to conditional-menu integrals as follows.
Let $\mathcal T=\{(t,\rho):0\le\rho\le t\le1\}$. Write
$r_c(t,\rho)$ and $n_c(t,\rho)$ for the expected payment and expected
item count of the single bidder facing $(t,\rho)$ under $M_*$, including
its lottery options, with its own two coordinates uniform. Let $P(t,\rho)$
and $B(t,\rho)$ be the scarce and safe singleton prices clipped at one.
With $g_R=P(1-P)+B(1-B)$ and $g_N=2-P-B$, the exact reductions are
\begin{equation}\label{eq:v5-face-reductions}
\begin{aligned}
 N_4&=4\int_{\mathcal T}n_c(t,\rho)\,dt\,d\rho,\\
 R_3&=\int_0^1r_c(t,0)\,dt+\int_{\mathcal T}g_R(t,\rho)\,dt\,d\rho,\\
 N_3&=\int_0^1n_c(t,0)\,dt+\int_{\mathcal T}g_N(t,\rho)\,dt\,d\rho,\\
 R_{\rm cross}&=2\int_0^1 B(t,0)(1-B(t,0))\,dt,\\
 N_{\rm cross}&=2\int_0^1(1-B(t,0))\,dt.
\end{aligned}
\end{equation}
On a one-zero face, the full two-coordinate bidder gives the line
integral, and the other bidder's singleton sales under the two opponent
orientations give $g_R$ and $g_N$. On a cross face only the safe singleton
threshold remains. Zero-coordinate exclusion removes any competing bundle
or lottery allocating the zero-valued item. The formulas and polynomial
branch integrands are also specified in
\path{verifier/primal_face_integrals.py} in the certificate package.

The pointwise tie rule need not be symmetric. The face integrals nevertheless have the required symmetries:
zero-coordinate exclusion handles fixed zero coordinates, ties involving
positive free coordinates lie on null own-type menu interfaces, and a
bidder whose two coordinates are fixed at zero selects empty. These facts
justify the face symmetries without changing the selected pointwise rule.

The verification partitions the opponent triangle into 296 fully resolved
strips and 15 wholly enclosed unresolved strips. Polynomial branch
conditions are certified over entire cells before exact integration;
unresolved strips receive valid global payment and allocation bounds.
The single-buyer radical is bracketed rationally, with a conservative
$10^{-12}$ face error allowance justified by the utility revenue identity.
An independent standard-library implementation reconstructs 1,372
polynomial cells and checks 63,184 exact Bernstein predicates before
integrating them. It agrees on every face enclosure and on the full
16-face expansion. It shares the analytic menu formulas and enclosure
principle, but imports no source calculator.

The same exact identity gives a global admissible first variation,
\[
 \left.\frac{d}{d\tau}R(M_\tau)\right|_{0+}
       =N_4+4R_3-5R_*
       \in[0.0120503811119742,0.0120521625636254].
\]
Thus this necessary optimality condition fails at $M_*$. The selected finite change at $\tau=83/10000$ gains more than $3/62500$
over the seed, as follows from the certified revenue and seed integral. Appendix~\ref{app:reserve-parameter} certifies the unique maximum within
this fixed-seed family and explains the selected rational parameter.
It does not identify an unrestricted optimal auction.
